\section*{第 \ref{ch:g12:phylogenetic-trees} 章——读懂亲缘关系：系统发生树}

\begin{solution}{exo:g12:phylogenetic-trees:1}
末端：所比较的\omterm{def:g10:biodiversity-scales:species}{物种}。节点：它们假想的最近
共同祖先。根：它们全部的共同祖先。
\end{solution}

\begin{solution}{exo:g12:phylogenetic-trees:2}
蜥蜴的最近亲缘类群是小鼠和人类（它们
共享蜥蜴分支正下方的那个节点）。青蛙的最近亲缘类群
是其余三者组成的整个类群，同样近。
\end{solution}

\begin{solution}{exo:g12:phylogenetic-trees:3}
一种在某一祖先中出现、并被其后代继承的状态；
共有它即意味着共有该祖先。祖先状态
为整个类群的祖先所具有，因此每个\omterm{def:g10:biodiversity-scales:species}{物种}都可能保留它：
共有它并不能说明更近的
亲缘关系。
\end{solution}

\begin{solution}{exo:g12:phylogenetic-trees:4}
一个祖先及其全部后代。鸟类是一个\omterm{prop:g12:phylogenetic-trees:rules}{演化支}。“爬行动物”在
日常意义上则不是：它遗漏了鸟类，而鸟类正是从其中
衍生出来的。
\end{solution}

\begin{solution}{exo:g12:phylogenetic-trees:5}
黑猩猩；最古老的节点
是连接猕猴与类人猿的那个节点，距今 2500 万年。
\end{solution}

\begin{solution}{exo:g12:phylogenetic-trees:6}
在每个节点处旋转：根先分为青蛙（底部）和其余类群；
再分出蜥蜴，最后是顶部的小鼠和人类。小鼠与人类仍然
共享最近的节点；蜥蜴接下来与它们相连，
青蛙最后加入。
\end{solution}

\begin{solution}{exo:g12:phylogenetic-trees:7}
鸽子具有颌、四肢和羊膜：它进入羊膜动物\omterm{prop:g12:phylogenetic-trees:rules}{演化支}，
与蜥蜴和小鼠并列。羽毛仅存在于一个\omterm{def:g10:biodiversity-scales:species}{物种}中，
在此不能界定任何类群；若加入更多鸟类，它们
便可界定鸟类。
\end{solution}

\begin{solution}{exo:g12:phylogenetic-trees:8}
毛发、乳汁和颌使海豚归入哺乳动物；若将它与鲨鱼归为一类，
则需要所有哺乳动物性状都
独立出现两次。鳍与流线型体形在两个谱系中
各自独立出现，是趋同演化。
\end{solution}

\begin{solution}{exo:g12:phylogenetic-trees:9}
$2.2/0.27 \approx 8$ 百万年。
\end{solution}

\begin{solution}{exo:g12:phylogenetic-trees:10}
二者均为现存末端，彼此都不是对方的后代。正确的说法是：
人类与黑猩猩是姐妹\omterm{def:g10:biodiversity-scales:species}{物种}，二者由一个共同祖先
衍生而来，该祖先已灭绝，生活于约 600 万年前，
既不是人类，也不是黑猩猩。
\end{solution}

\begin{solution}{exo:g12:phylogenetic-trees:11}
已知七鳃鳗位于有颌\omterm{def:g10:common-ancestry:bodyplan}{脊椎动物}类群之外，
因此其各状态对有颌\omterm{def:g10:common-ancestry:bodyplan}{脊椎动物}而言都是祖先状态。若以小鼠作外类群，
毛发、羊膜、四肢和颌都会被视为祖先状态，其缺失
则被视为衍生状态：系统树将整个倒置。
\end{solution}

\begin{solution}{exo:g12:phylogenetic-trees:12}
祖先\omterm{def:g12:selection-drift-speciation:population}{种群}的每个\omterm{def:g10:universal-dna:gene}{基因}都带有多个\omterm{def:g10:universal-dna:gene}{等位基因}；当它
在短时间内连续分化两次时，不同的\omterm{def:g10:universal-dna:gene}{等位基因}被分配
到三个后代中，且在不同\omterm{def:g10:universal-dna:gene}{基因}上的分配各不相同，因此某一\omterm{def:g10:universal-dna:gene}{基因}的历史
将\omterm{def:g10:biodiversity-scales:species}{物种} 1 与 2 归为一组，另一\omterm{def:g10:universal-dna:gene}{基因}的历史则将 2 与 3 归为一组。每个\omterm{def:g10:universal-dna:gene}{基因}树
对该\omterm{def:g10:universal-dna:gene}{基因}而言是正确的；\omterm{def:g10:biodiversity-scales:species}{物种}树则是众多\omterm{def:g10:universal-dna:gene}{基因}
的多数意见。
\end{solution}

\begin{solution}{exo:g12:phylogenetic-trees:13}
以每年 1\% 计，相隔数月采集的病毒样本已经
存在可测的差异：病毒树能够分辨一次流行
中谁感染了谁，但在数百万年间，序列会
被多次改写。以每百万年 0.3\% 计，哺乳动物的\omterm{def:g10:universal-dna:gene}{基因}
变化太慢，在一次流行中看不出任何变化，
但能在 2 亿年间保留可读的记录。
\end{solution}

\begin{solution}{exo:g12:phylogenetic-trees:14}
化石本身是一个独立的\omterm{def:g10:biodiversity-scales:species}{物种}，具有自己的衍生状态，
未必是祖先\omterm{def:g12:selection-drift-speciation:population}{种群}；它位于靠近羽毛出现节点的一条侧支上。
它表明该节点——鸟类与其最近的爬行动物亲缘类群的
共同祖先——距今至少 1.5 亿年，
并且羽毛先于牙齿与
尾巴的丧失而出现。
\end{solution}

\begin{solution}{exo:g12:phylogenetic-trees:15}
系统树没有顶端：旋转任一节点，人类就会移到底部，
而亲缘关系毫不改变。现存的每个末端自根部以来演化的时间
相同，因此没有哪一个更为高等。人类只是
数百万个末端中的一个，位于类人猿的一个分支上；系统树的形状记录的是
衍生关系，而不是等级。
\end{solution}

\begin{solution}{pb:g12:phylogenetic-trees:1}
\textbf{1.} 硬骨骼（五个\omterm{def:g10:biodiversity-scales:species}{物种}）：硬骨\omterm{def:g10:common-ancestry:bodyplan}{脊椎动物}；
鲨鱼不在其内。

\textbf{2.} 四肢：四足动物（蝾螈、鳄鱼、鸽子，
猫）。羊膜卵：羊膜动物（鳄鱼、鸽子、猫）。嵌套：
硬骨\omterm{def:g10:common-ancestry:bodyplan}{脊椎动物} $\supset$ 四足动物 $\supset$ 羊膜动物。

\textbf{3.} 否：仅见于单一\omterm{def:g10:biodiversity-scales:species}{物种}的状态无法划分类群。它们
若被更多\omterm{def:g10:biodiversity-scales:species}{物种}共有则有用（其他鸟类、其他
哺乳动物、其他爬行动物）。

\textbf{4.} 根：鲨鱼与其余（硬骨骼）；然后是鲑鱼与四足动物（四肢）；
然后是蝾螈与羊膜动物（羊膜卵）；
羊膜动物之中是一个未定的三分支，
鸽子一支具羽毛，猫的一支具毛发和
颅骨开口。

\textbf{5.} 表中将鳄鱼、鸽子和猫放在一起，
但无法为它们排序：鳄鱼最近的亲缘类群尚未确定。

\textbf{6.} 鳄鱼–鸽子，8\%。

\textbf{7.} 猫（与各者相差 15\%）；然后是蝾螈（19--20\%）；
然后是鲑鱼（24--25\%）；然后是鲨鱼（30--31\%）。

\textbf{8.} 二者在每一层嵌套上一致；\omterm{def:g10:universal-dna:information}{DNA}
分辨出羊膜动物：鳄鱼与鸽子互为姐妹群，猫随后并入。

\textbf{9.} 鲨鱼的谱系先于其余所有类群彼此分化而分开；
它们各自与鲨鱼
已分开相同的时间，因而所示距离相同。

\textbf{10.} 鳄鱼与鸽子的亲缘关系比与猫更近，
——按同样的道理，也比与蜥蜴更近；不含鸟类的“爬行动物”
不是一个\omterm{prop:g12:phylogenetic-trees:rules}{演化支}。

\textbf{11.} 19\%，分布于 $2 \times 340$ 百万年的分支上：约
每个谱系每百万年 0.028\%，即
每百万年分离 0.056\%。

\textbf{12.} 鳄鱼–鸽子 $8/0.056 \approx 140$ 百万年；
羊膜动物节点 $15/0.056 \approx 270$ 百万年。

\textbf{13.} 鲑鱼 $24.5/0.056 \approx 440$；鲨鱼 $30.5/0.056
\approx 540$ 百万年。

\textbf{14.} 化石：250 和 310；分子钟给出 140 和 270。
顺序正确，较老的那个日期也接近，但鳄鱼–鸟类的日期
明显过小：单个\omterm{def:g10:universal-dna:gene}{基因}的速率在不同谱系间有所不同，
单一\omterm{def:g10:universal-dna:gene}{基因}至多只能给出粗略的分子钟。

\textbf{15.} 距离较大时，许多位点已不止一次改变，因此差异
不再累加：分子钟趋于饱和。这一\omterm{def:g10:universal-dna:gene}{基因}低估了古老节点，
需要更慢的\omterm{def:g10:universal-dna:gene}{基因}，或是多个\omterm{def:g10:universal-dna:gene}{基因}，
方能胜任。

\textbf{16.} （猫，（鳄鱼，鸽子））是包含猫的最小\omterm{prop:g12:phylogenetic-trees:rules}{演化支}：
羊膜动物；其次是四足动物；再次是硬骨\omterm{def:g10:common-ancestry:bodyplan}{脊椎动物}；
最后是连同鲨鱼在内的全部六个\omterm{def:g10:biodiversity-scales:species}{物种}。

\textbf{17.} “鱼类”不是一个\omterm{prop:g12:phylogenetic-trees:rules}{演化支}：连接鲨鱼与鲑鱼的节点
即是根，其后代包括一切。羊膜动物是一个\omterm{prop:g12:phylogenetic-trees:rules}{演化支}：
鳄鱼、鸽子、猫及其共同祖先。

\textbf{18.} 蝾螈是现存\omterm{def:g10:biodiversity-scales:species}{物种}，而非中间类型；
它是羊膜动物的姐妹群，与之共享四足动物祖先，
但缺少羊膜卵，后者是在它们的谱系中较晚出现的。

\textbf{19.} 位于四足动物节点与羊膜动物节点之间分支的旁边：
一种不是羊膜动物的四足动物。其年代证实，
四肢在 360 百万年前即已存在，而羊膜卵
出现得更晚。

\textbf{20.} （鲨鱼，（鲑鱼，（蝾螈，（猫，（鳄鱼，
鸽子）））））；\omterm{def:g10:universal-dna:information}{DNA} 确定鳄鱼与鸽子互为姐妹群；
依据该\omterm{def:g10:universal-dna:gene}{基因}，二者分化于约 140 百万年前（化石
记录为 250 百万年）。
\end{solution}
